\VerdierPhysicalPageAnchor{0076}%
\VerdierSemanticAnchor{ch1-2-4-2-countable-products-sums}
Suppose now that countable products (respectively, direct sums) are
representable in $C$. Set:

\VerdierSemanticAnchor{ch1-2-4-2-eq6}
\begin{equation}
  \begin{aligned}
    \operatorname{comp}_{\Pi}F
      &=\int_n^{\Pi}F^{[n]},\\
    \text{(respectively, }\operatorname{comp}_{\Sigma}F
      &=\int_n^{\Sigma}F^{[n]}\text{)}.
  \end{aligned}
  \tag{2.4.2.6}
\end{equation}

\VerdierSemanticAnchor{ch1-2-4-2-pi-sigma-extension}
The functor
$\operatorname{comp}_{\Pi}F$ (respectively,
$\operatorname{comp}_{\Sigma}F$) :
$\prod_{i\in[n]}\operatorname{comp}(C_i)\longrightarrow
\operatorname{comp}(C)$ is called the $\Pi$-extension (respectively, the
$\Sigma$-extension) of the functor $F$ to complexes.

\VerdierSemanticAnchor{ch1-2-4-2-koszul-twist}
Since $F$ is a strict $\mathbb Z[n]$-functor and the functor
$\int_n^{\Pi}$ (respectively, $\int_n^{\Sigma}$) is a
$\mathbb Z[n]$-functor twisted by the Koszul cocycle relative to the ordered
basis $e_1,\ldots,e_n$ of $\mathbb Z[n]$, the functor
$\operatorname{comp}_{\Pi}F$ (respectively,
$\operatorname{comp}_{\Sigma}F$) is a $\mathbb Z[n]$-functor twisted by the
Koszul cocycle relative to the ordered basis $e_1,\ldots,e_n$ of
$\mathbb Z[n]$, in accordance with Sign Rule I (1.6.6). Thus, for every
$i\in[n]$, one has canonical isomorphisms:

\VerdierSemanticAnchor{ch1-2-4-2-eq7}
\begin{equation}
  \begin{aligned}
    \operatorname{comp}_{\Pi}F(X_1,\ldots,X_i[1],\ldots,X_n)
      &\xrightarrow{\sim}
      \operatorname{comp}_{\Pi}F(X_1,\ldots,X_i,\ldots,X_n)[1],\\
    \text{(respectively, }\operatorname{comp}_{\Sigma}F(X_1,\ldots,X_i[1],\ldots,X_n)
      &\xrightarrow{\sim}
      \operatorname{comp}_{\Sigma}F(X_1,\ldots,X_i,\ldots,X_n)[1]\text{)}
  \end{aligned}
  \tag{2.4.2.7}
\end{equation}

\VerdierSemanticAnchor{ch1-2-4-2-canonical-pi}
which will be denoted by $p_i$.

\VerdierSemanticAnchor{ch1-2-4-2-anticommutativity}
These isomorphisms have the anticommutativity properties described in
(1.6.6.3).

\VerdierSemanticAnchor{ch1-2-4-2-summable-restriction-equality}
The restrictions of the functors $\operatorname{comp}_{\Pi}F$ and
$\operatorname{comp}_{\Sigma}F$ to the summable objects of
$\prod_{i\in[n]}\operatorname{comp}(C_i)$ are equal, and this restriction
can always be defined without hypotheses on the category $C$.

\VerdierSemanticAnchor{ch1-2-4-3}
\subsection*{2.4.3.}
Now let
$F,F'\colon\prod_{i\in[n]}C_i\longrightarrow C$ be two multiadditive
functors, and let $h\colon F\longrightarrow F'$ be a morphism of functors.
From it one obtains a morphism of functors:

\VerdierSemanticAnchor{ch1-2-4-3-induced-morphism}
\begin{equation*}
  h^{[n]}\colon F^{[n]}\longrightarrow F'^{[n]},
\end{equation*}
by setting:

\VerdierSemanticAnchor{ch1-2-4-3-eq1}
\begin{equation}
  \begin{aligned}
    h^{[n]}(X_1,\ldots,X_n)^{\sigma}
      &=h(X_1^{l_1},\ldots,X_n^{l_n}),\\
    \sigma&=\sum_{i\in[n]}l_i e_i\in\mathbb Z[n],\\
    (X_1,\ldots,X_n)&\in
      \operatorname{Ob}\!\left(\prod_{i\in[n]}\operatorname{comp}(C_i)\right).
  \end{aligned}
  \tag{2.4.3.1}
\end{equation}
